\documentclass[letterpaper,10pt]{texMemo}

\usepackage{graphicx}

\memoto{Dr. Mark Yoder}
\memofrom{Michael McDonald}
\memosubject{Week 3 (I$^2$C, PWM, Device Tree Overlays) Memo}
\memodate{\today}
\logo{\includegraphics[width=0.5\textwidth]{RoseLogo.png}}

\begin{document}
\maketitle

\section{Exercise Summary}

\noindent I have completed Homework 3.

\begin{center}
\begin{tabular}{|p{4cm}|l|l|p{8cm}|}
\hline
\textbf{Objective} & \textbf{Exercise} &\textbf{Status} & \textbf{Notes}\\
\hline
Device Tree Overlays & Exercise 11a & Completed & These pins seem to work, but
I'm worried that the pins we need to use for our homework won't since they
aren't free.\\
\hline
I$^2$C Communication & Exercise 12 & Completed & The only thing I noticed is
that myi2cget requires the decimal versions of addresses, rather than hex.\\
\hline
PWM Control & Exercise 13 & Completed & A 1MHz signal is very crisp, but the
10MHz signal was exactly as described (~9MHz and not clean at all). There's also
a pretty impressive (~1V) overshoot above 3.3V for the faster signals.\\
\hline
\end{tabular}
\end{center}

\noindent The main points I took away from these exercises were the following:
\begin{itemize}
 \item The on-board hardware is awesome (both I$^2$C and PWM). It makes life so
much easier, and decreases development time (and reduces frustration). That
being said, it's much harder to understand how all of the layers interact
because there are so many more layers between the software and the hardware.
 \item There's a difference between decimal and hexidecimal, and it can come
around to bite you when doing I$^2$C if you aren't careful about which
functions take what as an argument.
\end{itemize}

\end{document}
